\chapter[PENDAHULUAN]{PENDAHULUAN}
\pagenumbering{arabic}
\sloppy

\section{Latar Belakang}

Sektor industri ritel fisik (\textit{brick-and-mortar}) terus beradaptasi mempertahankan daya saing operasional di tengah pesatnya pertumbuhan ekosistem perdagangan digital. Meskipun mayoritas transaksi belanja masyarakat masih bertumpu pada kehadiran fisik di toko, para pengelola ritel menghadapi tantangan operasional yang kian berat, mulai dari lonjakan biaya tenaga kerja, ketatnya kompetisi ruang usaha, hingga tuntutan peningkatan efisiensi layanan secara berkesinambungan. Untuk menjawab tantangan tersebut, sektor ritel modern mengalami pergeseran paradigma menuju otomasi toko pintar (\textit{smart stores}). Salah satu pilar penentu dalam transformasi ini adalah kemampuan pengelola untuk mengekstraksi data perilaku konsumen secara akurat, terukur, dan objektif sebagai landasan pengambilan keputusan strategis.

Secara konvensional, pemahaman mengenai perilaku pengunjung di dalam toko ritel sangat bergantung pada metode manual, seperti observasi petugas, pencatatan kasir, maupun survei kepuasan pelanggan secara berkala. Namun, pendekatan konvensional tersebut memiliki kelemahan mendasar karena membutuhkan alokasi sumber daya manusia yang besar, rentan terhadap bias subjektif, serta tidak mampu menyediakan data operasional secara berkesinambungan. Penerapan teknologi sensor fisik seperti kartu identifikasi RFID atau gerbang pemindai inframerah juga menghadapi hambatan tersendiri, baik dari segi biaya investasi infrastruktur yang tinggi maupun keengganan pelanggan untuk berinteraksi dengan perangkat pemindai yang invasif. Keterbatasan instrumen tradisional tersebut mendorong urgensi adopsi teknologi visi komputer (\textit{computer vision}) yang mampu mengolah rekaman video pengawas (CCTV) secara non-invasif untuk memantau dinamika konsumen secara pasif \cite{fierrosilva2026, hairurizal2026}.

Integrasi visi komputer pada ekosistem komersial memungkinkan sistem analitik video untuk mengekstraksi berbagai indikator kinerja utama secara otomatis dan waktu nyata. Data visual yang dihimpun dari kamera pengawas dapat dikonversi menjadi metrik perilaku yang sangat bernilai bagi pengelola toko, seperti volume kunjungan (\textit{people counting}), tingkat kepadatan area (\textit{occupancy}), persebaran atensi pengunjung terhadap etalase produk, hingga durasi waktu singgah (\textit{dwell time}) pelanggan di depan rak pajang \cite{shili2024electronics, narvilas2022}. Pemahaman spasial dan temporal mengenai rute penjelajahan konsumen tersebut menjadi instrumen krusial dalam mengevaluasi efektivitas tata letak gerai, mengidentifikasi produk yang paling diminati, hingga mengantisipasi penumpukan antrean transaksi pada meja pembayaran.

Akan tetapi, implementasi sistem analitik visual pada denah toko ritel skala penuh tidak dapat bertumpu hanya pada penggunaan kamera tunggal (\textit{single-camera}). Kamera tunggal memiliki keterbatasan sudut pandang (\textit{field-of-view}) serta menghadirkan area titik buta (\textit{blind spots}) yang cukup luas akibat tertutup oleh partisi lorong belanja atau etalase bertingkat. Ketika seorang konsumen melangkah keluar dari cakupan satu lensa kamera, sistem pelacakan otomatis kehilangan jejak individu tersebut. Kondisi ini menuntut penerapan konfigurasi jaringan multi-kamera (\textit{multi-camera system}) guna memperluas cakupan area pemantauan secara holistik melintasi berbagai zona yang tidak saling bertumpukan (\textit{non-overlapping views}) \cite{fierrosilva2026}. Melalui sinkronisasi beberapa sudut pandang kamera, pergerakan pelanggan diharapkan dapat terus dipantau secara berkesinambungan dari pintu masuk hingga area pembayaran.

Kendati konfigurasi multi-kamera menawarkan cakupan ruang yang komprehensif, pelacakan pergerakan manusia di lingkungan komersial dunia nyata menghadirkan tantangan teknis yang sangat dinamis. Pola pergerakan konsumen di dalam toko sangat acak dan tidak terduga; pengunjung kerap berhenti mendadak, berputar arah, atau berkumpul dalam kelompok padat di area promosi tertentu. Kondisi kepadatan ini memicu timbulnya fenomena oklusi visual, yakni situasi di mana tubuh seorang pengunjung tertutup sebagian maupun seluruhnya oleh struktur rak barang, troli belanja, atau oleh pengunjung lain yang melintas di depannya \cite{panagos2026, mendes2023multi}. Pada situasi oklusi berat, pandangan kamera terhadap target pengamatan terputus sementara waktu, sehingga sistem analitik kehilangan kesinambungan visual atas subjek yang sedang diamati.

Dampak paling mendasar yang ditimbulkan oleh oklusi visual adalah terjadinya fragmentasi identitas objek atau fenomena \textit{identity churn} \cite{fierrosilva2026}. Ketika seorang pelanggan yang terhalang selama beberapa detik kembali terlihat oleh kamera, sistem pelacakan kerap gagal mengenali bahwa orang tersebut adalah subjek yang sama dengan yang terpantau sebelumnya. Alih-alih menyambung riwayat lintasan yang sudah ada, algoritma pelacak justru menetapkan nomor identitas (ID) baru kepada pelanggan tersebut. Fenomena pemutusan dan pergantian ID ini berdampak fatal terhadap keabsahan penghitungan metrik analitik bisnis. Sebagai ilustrasi, seorang pelanggan yang sebenarnya mengantre di kasir selama lima menit dapat terhitung secara keliru sebagai tiga orang berbeda yang masing-masing hanya berdiri selama satu koma lima menit, akibat tubuhnya sempat tertutup oleh pengunjung lain yang melintas. Hal ini menghasilkan laporan analitik yang salah kaprah dan merugikan pihak manajemen dalam menilai kualitas layanan antrean.

Selama ini, perkembangan literatur di bidang visi komputer dan pelacakan multi-objek (\textit{Multi-Object Tracking}) didominasi oleh upaya peningkatan performa algoritma pada metrik evaluasi standar laboratorium, seperti \textit{Multiple Object Tracking Accuracy} (MOTA) dan \textit{Identification F1-Score} (IDF1) \cite{bayraktar2025retrackvlm}. Namun, terdapat kesenjangan pengetahuan yang cukup lebar antara pencapaian teoretis di laboratorium dengan kebutuhan implementasi di lapangan. Pengujian akademik umumnya dilakukan pada dataset publik yang terisolasi dengan skenario gerak yang relatif teratur. Sangat jarang ditemukan penelitian empiris yang menguji sejauh mana keunggulan skor metrik pelacakan akademik tersebut berbanding lurus dengan ketepatan penghitungan metrik bisnis perilaku konsumen saat dihadapkan pada skenario ritel komersial yang penuh dengan oklusi dan interaksi padat. Sistem yang mencatatkan nilai akurasi pelacakan tinggi secara matematis belum tentu menjamin ketepatan penghitungan durasi \textit{dwell time} atau waktu antre jika algoritma tersebut rentan mengalami fragmentasi identitas.

Berangkat dari kesenjangan tersebut, penelitian tugas akhir ini diarahkan untuk melakukan evaluasi komparatif secara objektif terhadap performa berbagai algoritma pelacakan multi-objek mutakhir dalam mengukur indikator perilaku konsumen pada skenario ritel fisik. Pengujian difokuskan pada ketahanan algoritma pelacakan dalam mempertahankan keutuhan identitas target di tengah kondisi oklusi yang kerap terjadi di area perbelanjaan. Melalui perbandingan langsung antara hasil keluaran sistem analitik otomatis dengan data observasi manual di lapangan (\textit{ground truth}), penelitian ini bertujuan membuktikan relevansi korelasi antara metrik tolok ukur akademik dengan presisi metrik bisnis ritel, sekaligus mengidentifikasi pendekatan pelacakan yang paling tangguh untuk meminimalkan margin kesalahan penghitungan di lingkungan operasional nyata.


\section{Rumusan Masalah}

Berdasarkan latar belakang yang telah dipaparkan, rumusan masalah dalam penelitian tugas akhir ini adalah sebagai berikut:

\begin{enumerate}
    \item Bagaimana tingkat akurasi pengukuran metrik perilaku konsumen, yang mencakup durasi waktu singgah (\textit{dwell time}), waktu tunggu antrean, dan estimasi jumlah pengunjung (\textit{counting}), menggunakan sistem analitik berbasis \textit{Multi-Object Tracking} pada kondisi lingkungan ritel nyata dengan tingkat kepadatan dan oklusi tinggi?
    \item Algoritma pelacakan manakah di antara metode kontemporer (ByteTrack, BoT-SORT, OC-SORT, dan Deep OC-SORT) yang menghasilkan tingkat galat metrik perilaku konsumen terkecil bila dibandingkan dengan data aktual di lapangan (\textit{ground truth}) maupun platform analitik video komersial?
    \item Bagaimana derajat korelasi empiris antara capaian metrik evaluasi pelacakan standar akademik (seperti MOTA dan IDF1) dengan tingkat presisi pengukuran metrik perilaku konsumen di sektor komersial dunia nyata?
\end{enumerate}

\section{Batasan Masalah atau Ruang Lingkup}

Batasan masalah dalam penelitian tugas akhir ini ditetapkan sebagai berikut:

\begin{enumerate}
    \item Analisis pemantauan visual hanya difokuskan pada deteksi target manusia (kategori kelas \textit{person}).
    \item Arsitektur model pendeteksi objek dikunci menggunakan YOLO11 (satu versi model dan bobot tetap) sebagai variabel kontrol, sehingga variasi hasil pengujian murni mencerminkan performa algoritma pelacakan yang dievaluasi.
    \item Evaluasi algoritma pelacakan multi-objek dibatasi pada empat metode mutakhir: ByteTrack, BoT-SORT, OC-SORT, dan Deep OC-SORT.
    \item Pelacakan identitas konsumen ditargetkan pada interval kunjungan lokal jangka pendek (skala detik hingga menit antar-bingkai video), tanpa melibatkan re-identifikasi jangka panjang lintas hari.
    \item Validasi tingkat galat sistem dievaluasi terhadap data anotasi manual (\textit{ground truth}) dari rekaman video CCTV berdurasi 10--15 menit pada jam operasional sibuk di dua lokasi area komersial.
    \item Pemrosesan data video dilakukan secara anonim demi menjaga privasi pengunjung tanpa menerapkan modul pengenalan wajah (\textit{facial recognition}).
\end{enumerate}

\section{Tujuan}

Sejalan dengan rumusan masalah yang diajukan, tujuan dari pelaksanaan penelitian tugas akhir ini adalah sebagai berikut:

\begin{enumerate}
    \item Mengukur secara kuantitatif tingkat akurasi penghitungan metrik perilaku konsumen (durasi \textit{dwell time}, waktu tunggu antrean, dan estimasi jumlah pengunjung) yang diproses oleh sistem analitik berbasis \textit{Multi-Object Tracking} di bawah pengaruh gangguan oklusi pada area ritel.
    \item Mengevaluasi dan menentukan algoritma pelacakan multi-objek kontemporer yang menghasilkan nilai galat terendah terhadap pengamatan manual di lapangan (\textit{ground truth}) serta membandingkannya secara objektif dengan performa platform analitik video komersial.
    \item Menganalisis derajat korelasi empiris antara capaian metrik evaluasi pelacakan standar akademik (MOTA dan IDF1) dengan tingkat validitas data analitik perilaku konsumen di lapangan guna membuktikan relevansi tolok ukur teoretis terhadap kebutuhan praktis industri.
\end{enumerate}

\section{Manfaat}

Manfaat yang diharapkan dari pelaksanaan penelitian ini adalah sebagai berikut:

\begin{enumerate}
    \item Memberikan kontribusi teoretis dengan mengisi celah literatur terkait validitas penggunaan metrik evaluasi pelacakan konvensional (MOTA dan IDF1) sebagai prediktor keandalan metrik analitik bisnis di dunia nyata.
    \item Menyediakan rekomendasi teknis berbasis bukti empiris bagi pengembang perangkat lunak maupun pelaku industri ritel dalam memilih algoritma pelacakan objek yang paling andal terhadap kendala oklusi visual dan fragmentasi identitas untuk diterapkan pada arsitektur platform analitik video pintar.
\end{enumerate}